\documentclass[11pt]{article}

\usepackage[utf8]{inputenc}
\usepackage[T1]{fontenc}
\usepackage[margin=1in]{geometry}
\usepackage{enumitem}
\usepackage[colorlinks=true,linkcolor=blue,linktoc=all]{hyperref}

% ==========================================================
% xr-hyper pulls in every \label from the constitution's own
% main.aux, so every \ref{...} below resolves to whatever Article/
% Section number that label currently has in the constitution --
% automatically, even if the constitution is reorganized again
% later. This document must be compiled from the SAME DIRECTORY
% as the constitution's main.aux (or the path below adjusted),
% and the constitution must be compiled at least once first so
% main.aux exists and is current.
% ==========================================================
\usepackage{xr-hyper}
\externaldocument[const-]{main}

\title{Sortition Legislature: Operational Responsibilities Guide}
\author{Companion document to \emph{A Constitution for the Commonwealth} -- not part of the constitutional text}
\date{}

\begin{document}
\maketitle

\begin{quote}
\itshape
This guide is an operational aid for members of the Sortition Legislature. It has no constitutional force of its own: where anything here appears to conflict with the constitutional text, the constitution controls. Every citation below links to the specific Article and Section that is the actual source of authority.
\end{quote}

\tableofcontents

\section{How to Use This Guide}

The Legislature's constitutional responsibilities fall into four categories, each with a different operational logic:

\begin{enumerate}
  \item \textbf{Standing Duties} -- always in force; no trigger or schedule activates them, they simply describe how the Legislature conducts itself at all times.
  \item \textbf{Cyclical Reviews} -- recur automatically on a fixed calendar, regardless of any external event.
  \item \textbf{Deadline-Triggered Actions} -- activate only when a specific event occurs (a threshold crossing, a report, a petition), and carry a hard deadline once triggered. \textbf{These are the highest-priority items in this guide}, since missing a constitutional deadline has direct legal consequences (e.g., a provisional prohibition becoming permanent).
  \item \textbf{Discretionary Powers} -- the Legislature may exercise these at any time, at its own initiative, with no deadline forcing action.
\end{enumerate}

Within Section~3 (Deadline-Triggered Actions), entries are ordered by deadline length, shortest first, since that is the order in which they would actually compete for the Legislature's attention if several were active simultaneously.

\section{Standing Duties}

These are always in force. They describe the Legislature's ongoing conduct, not a discrete task to complete.

\begin{enumerate}
  \item \textbf{Follow the structured deliberative process for all significant legislation} -- expert briefings from multiple perspectives, small-group deliberation in randomly assigned groups, full-assembly reporting and cross-group deliberation, minority impact assessment, and final deliberation and decision, in that order. (Article~\ref{const-art:sortition-legislature}, \ref{const-art:sortition-legislature:s-deliberative-process})
  \item \textbf{Conduct all proceedings with real-time interpretation} in every language spoken by more than one percent of the Commonwealth population, and make official documents and civic materials available in all threshold languages. (Article~\ref{const-art:sortition-legislature}, \ref{const-art:sortition-legislature:s-multilingualism})
  \item \textbf{Maintain individual conflict-of-interest compliance} -- full financial and relational disclosure before service; no income beyond public salary and civic service premium; no private communication with entities interested in pending decisions; no productive assets beyond personal property. This applies to all governance roles, not only the Legislature. (Article~\ref{const-art:governance-principles}, \ref{const-art:governance-principles:s-conflict-interest-sector})
  \item \textbf{Observe sector exclusion} -- no member may deliberate on a matter in which they stand to gain or lose materially in a way distinguishable from the general population. This applies to all governance roles, not only the Legislature. (Article~\ref{const-art:governance-principles}, \ref{const-art:governance-principles:s-conflict-interest-sector})
  \item \textbf{Fund the Material Sufficiency Floor on an ongoing basis} through cooperative surplus contribution rates set at a level the Monetary Authority and Financial Regulatory Body jointly assess as structurally sufficient. (Article~\ref{const-art:finance-money}, \ref{const-art:finance-money:s-cooperative-surplus-contributions})
\end{enumerate}

\section{Cyclical Reviews}

These recur automatically on a fixed calendar. No external trigger is needed -- the Legislature should track these dates proactively.

\begin{enumerate}
  \item \textbf{Climate mandate review} -- every five years, with automatic escalation triggers tied to objective scientific metrics. (Article~\ref{const-art:ecological-limits}, \ref{const-art:ecological-limits:s-climate-mandate})
  \item \textbf{Cooperative surplus contribution rate review} -- every five years, timed alongside the climate mandate review above. (Article~\ref{const-art:finance-money}, \ref{const-art:finance-money:s-cooperative-surplus-contributions})
  \item \textbf{Essential-sector forward-looking capacity targets and Incentive Escalation Schedules} -- reviewed every five years, alongside the same climate mandate cycle. (Article~\ref{const-art:essential-labor}, \ref{const-art:essential-labor:s-voluntary-service-incentives} and \ref{const-art:essential-labor:s-forward-looking-capacity})
  \item \textbf{Regulatory Body sector-exclusion thresholds} -- the time period defining a "recently adjudicated" matter, and the threshold for disqualifying outside income, are reviewed annually. (Article~\ref{const-art:regulatory-bodies}, \ref{const-art:regulatory-bodies:s-sector-exclusion})
\end{enumerate}

\section{Deadline-Triggered Actions}

Each entry below activates only when its trigger occurs. Ordered shortest deadline first.

\begin{enumerate}

  \item \textbf{Trigger: an emergency threshold is confirmed by the Emergency Threshold Verification Body.}
  \\ \textbf{Deadline: convene within 72 hours}, with a quorum of at least one third of members. The Legislature may authorize specific time-limited actions beyond the pre-approved protocol's scope; these authorizations expire after 30 days unless actively reauthorized. (Article~\ref{const-art:emergency-protocols}, \ref{const-art:emergency-protocols:s-novel-emergencies-three})

  \item \textbf{Trigger: a Legislative Decision Review petition reaches its civic-membership signature threshold (5\%).}
  \\ \textbf{Deadline: mandatory reconsideration session within 60 days} of threshold confirmation. Reconsideration does not suspend the original decision and does not bind the Legislature to reversal. (Article~\ref{const-art:civic-initiative-infrastructure}, \ref{const-art:civic-initiative-infrastructure:s-constitutional-intervention-types})

  \item \textbf{Trigger: the Monetary Authority issues a structural dependence finding} (allocation substituting for inadequate cooperative surplus contribution rates).
  \\ \textbf{Deadline: mandatory review session within 60 days} of receiving the report. (Article~\ref{const-art:finance-money}, \ref{const-art:finance-money:s-monetary-authority})

  \item \textbf{Trigger: the Transparency Engine issues an automatic alert} that Material Sufficiency Floor obligations are not structurally fully funded.
  \\ \textbf{Deadline: mandatory review session within 60 days.} (Article~\ref{const-art:finance-money}, \ref{const-art:finance-money:s-cooperative-surplus-contributions})

  \item \textbf{Trigger: an essential-sector Implementation Council reports a reactive enrollment threshold crossing (workforce pipeline risk).}
  \\ \textbf{Deadline: mandatory review session within 60 days} to assess and strengthen voluntary incentive mechanisms. (Article~\ref{const-art:essential-labor}, \ref{const-art:essential-labor:s-essential-sector-pipeline})

  \item \textbf{Trigger: an essential-sector Implementation Council reports a forward-looking capacity shortfall} (ten-year horizon), even before any reactive threshold is crossed.
  \\ \textbf{Deadline: mandatory review session within 60 days} to strengthen incentives, expand training pipeline capacity, or both. (Article~\ref{const-art:essential-labor}, \ref{const-art:essential-labor:s-forward-looking-capacity})

  \item \textbf{Trigger: emergency actions are authorized under the 72-hour emergency session (above).}
  \\ \textbf{Deadline: full Legislature reconvenes within 30 days} to review all actions taken under the emergency authorization; no extension without this active reauthorization. (Article~\ref{const-art:emergency-protocols}, \ref{const-art:emergency-protocols:s-novel-emergencies-three})

  \item \textbf{Trigger: a Regulatory Commons Body classifies a financial instrument as provisionally prohibited.}
  \\ \textbf{Deadline: the classification is referred to the Legislature within 30 days} of classification (a duty on the Financial Regulatory Body, not the Legislature); \textbf{the Legislature then has 90 days from referral to act}, or the provisional prohibition becomes permanent automatically. (Article~\ref{const-art:regulatory-bodies}, \ref{const-art:regulatory-bodies:s-financial-regulation})

\end{enumerate}

\section{Discretionary Powers}

The Legislature may exercise these at its own initiative. No deadline compels action, though some are the necessary predicate for other bodies' ability to function.

\subsection{Constitutional and Structural}
\begin{itemize}
  \item Amend constitutional provisions outside the Title~\ref{const-title:foundational} foundational principles, by four-fifths supermajority, followed by 60\% national ratification and Constitutional Court review. (Article~\ref{const-art:ordinary-amend})
  \item Ratify treaties negotiated by the Foreign Relations Council; the Council may negotiate but never conclude a binding agreement without this ratification. (Article~\ref{const-art:foreign-relations}, \ref{const-art:foreign-relations:s-mandate-constraints})
  \item Add new Civic Initiative Infrastructure intervention types (may not remove the constitutionally established ones). (Article~\ref{const-art:civic-initiative-infrastructure}, \ref{const-art:civic-initiative-infrastructure:s-constitutional-intervention-types})
  \item Create new Implementation Councils by supermajority vote when new governance domains emerge; dissolve existing ones by supermajority plus Constitutional Court confirmation that no Material Sufficiency Floor function is left unimplemented. (Article~\ref{const-art:implementation-councils}, \ref{const-art:implementation-councils:s-creation-dissolution})
  \item Create additional Regulatory Commons Bodies by supermajority vote (constitutionally mandated Bodies may only be dissolved via amendment). (Article~\ref{const-art:regulatory-bodies}, \ref{const-art:regulatory-bodies:s-financial-regulation})
  \item Define and adjust the mandate of each Regulatory Commons Body. (Article~\ref{const-art:regulatory-bodies}, \ref{const-art:regulatory-bodies:s-mandate-independence})
  \item Pre-approve detailed emergency response protocols during normal times, for each predictable emergency category, designating lead and supporting Implementation Councils in advance. (Article~\ref{const-art:emergency-protocols}, \ref{const-art:emergency-protocols:s-pre-legislation-requirement})
  \item Ratify Electoral Commons stratification-methodology changes and subsidiary exemption criteria. (Article~\ref{const-art:sortition-design}, \ref{const-art:sortition-design:s-stratification} and \ref{const-art:sortition-design:s-mandatory-service-civic})
\end{itemize}

\subsection{Economic}
\begin{itemize}
  \item Lower (never raise) the wage ratio cap on cooperative compensation. (Article~\ref{const-art:worker-coop-sector}, \ref{const-art:worker-coop-sector:s-wage-ratio-cap})
  \item Set or raise (never lower) essential-sector minimum compensation floors, both for cooperative-organized and Commons-institution essential labor. (Article~\ref{const-art:essential-labor}, \ref{const-art:essential-labor:s-essential-labor})
  \item Raise cooperative surplus contribution rates by simple majority; lower them only by supermajority with a published funding-impact assessment (and never while a structural-dependence finding from the preceding two years is outstanding). (Article~\ref{const-art:finance-money}, \ref{const-art:finance-money:s-cooperative-surplus-contributions})
  \item Raise essential-sector training stipends above the statutory floor, including in response to pipeline monitoring triggers. (Article~\ref{const-art:essential-labor}, \ref{const-art:essential-labor:s-training-pipelines})
  \item Expand incentive benefits beyond a published Incentive Escalation Schedule, or revise the Schedule itself for future shortfalls (may not reduce benefits already active for persons currently serving). (Article~\ref{const-art:essential-labor}, \ref{const-art:essential-labor:s-voluntary-service-incentives})
  \item Lift any provisional or permanent financial prohibition by simple majority vote (the Legislature alone may create new categories of permitted financial activity; the Financial Regulatory Body may not). (Article~\ref{const-art:regulatory-bodies}, \ref{const-art:regulatory-bodies:s-financial-regulation})
\end{itemize}

\end{document}